% Protocols are systems: blueprint.
% Core results live in systems-science-foundations; the three protocol
% instances live in this repository under Protocols/.

\chapter{Systems and the floor}

\begin{definition}[System]
  \label{def:klir-system}
  \lean{Systems.KlirSystem}
  \leanok
  A system is a pair $S = (T, R)$ of a set of things $T$ and a relation
  $R \subseteq T \times T$.
\end{definition}

\begin{remark}
  Klir, \emph{Facets of Systems Science} (2001), p.~5: ``S, T, R denote,
  respectively, a system, a set of things distinguished within S, and a
  relation (or, possibly, a set of relations) defined on T. Clearly, the
  thinghood and systemhood properties of S reside in T and R, respectively.''
  His examples on the same page admit both an ordering, which runs one way,
  and a partition, which comes from an equivalence relation and runs both ways.
\end{remark}

\begin{definition}[Shape of a definition]
  \label{def:klir-shape}
  \lean{KlirPosition, KlirArrow, KlirShape}
  \leanok
  \uses{def:klir-system}
  The shape $\IKlir$ has two objects, the relation and the things, and one
  generating arrow
  \[
    R \longrightarrow T,
  \]
  read ``$R$ is defined over $T$''. The arrow joins the parts of the
  definition, not two things.
\end{definition}

\begin{definition}[Kernel]
  \label{def:kernel}
  \lean{Systems.Kernel}
  \leanok
  \uses{def:klir-system, def:klir-shape}
  A kernel is a system $(T, R)$ together with the constraint that every pair
  in $R$ lies in $T$. The constraint is the arrow of
  Definition~\ref{def:klir-shape} at the level of data.
\end{definition}

\begin{theorem}[The arrow is in every definition]
  \label{thm:existence}
  \lean{klirToBunge_obj_injective, klirToMobus_obj_injective,
    klirToMyers_obj_injective, klirToWymore_obj_injective,
    klirToMesarovic_obj_injective, klirToJoslyn_obj_injective,
    klirToSpivak_obj_injective}
  \leanok
  \uses{def:klir-shape}
  For each of the definitions of Bunge, Mobus, Myers, Wymore, Mesarovi\'c,
  Joslyn, and Spivak, there is a functor from $\IKlir$ to that definition's
  shape that is injective on objects. With Klir's own, the arrow sits inside
  all eight.
\end{theorem}

\begin{proof}
  \leanok
  Each functor is given by a choice of positions in the target shape.
  Faithfulness is automatic because $\IKlir$ is thin, so the content is
  injectivity on objects.
\end{proof}

\begin{lemma}[No two arrows in a row]
  \label{lem:no-two-chain}
  \lean{SharedPrimitive.no_two_chain}
  \leanok
  \uses{def:klir-shape}
  No quiver that embeds in the Willems shape contains a path of two arrows.
\end{lemma}

\begin{proof}
  \leanok
  The Willems shape has no composable pair of arrows.
\end{proof}

\begin{theorem}[Only one shared arrow]
  \label{thm:single-arrow}
  \lean{SharedPrimitive.connected_is_single_arrow}
  \leanok
  \uses{def:klir-shape, lem:no-two-chain}
  Let $V$ be a connected quiver that embeds in both the Joslyn and the
  Willems shapes, and let $e : x \to y$ be an edge of $V$. Then $x \neq y$,
  $V = \{x, y\}$, and every edge of $V$ runs from $x$ to $y$.
\end{theorem}

\begin{proof}
  \leanok
  Joslyn's shape forbids two arrows leaving one object, and Willems's forbids
  two arrows entering one object and any two arrows in a row
  (Lemma~\ref{lem:no-two-chain}). A connected quiver satisfying all three
  has one edge.
\end{proof}

\begin{remark}
  The result is relative to the encoded presentations. Willems serves as a
  witness and is not one of the eight. Among things, a relation may run both
  ways. Among the parts of a definition, the shared floor is one arrow.
\end{remark}

\begin{theorem}[Nothing larger is shared, refuted]
  \label{thm:maximality-fails}
  \lean{SharedPrimitive.free_category_maximality_fails}
  \leanok
  \uses{def:klir-shape}
  Compared as free categories, the three-object fork embeds injectively on
  objects and faithfully into all eight shapes. So the claim that nothing
  larger than one arrow is shared is false at that level.
\end{theorem}

\begin{proof}
  \leanok
  In a free category every path is a morphism, and the fork enters through
  composite paths that no definition asserts.
\end{proof}

\chapter{What the floor can and cannot carry}

\begin{theorem}[Views are generated from the kernel]
  \label{thm:view-roundtrip}
  \lean{Systems.Kernel.toBunge_toKlir, Systems.Kernel.toMobus_toKlir}
  \leanok
  \uses{def:kernel}
  A kernel with at least one bonded pair of distinct things generates a Bunge
  system, and a kernel in which nothing is related to itself generates a Mobus
  system. Each generated system projects back to the same kernel.
\end{theorem}

\begin{proof}
  \leanok
  The elaboration slots are filled with minimal witnesses: empty environment,
  no external flows, trivial boundary, history, and timescale. The relation
  passes through unchanged.
\end{proof}

\begin{theorem}[The generated energy view is static]
  \label{thm:spivak-static}
  \lean{Systems.Kernel.toSpivak_static}
  \leanok
  \uses{def:kernel}
  The Spivak system generated from any kernel is static. Motion needs value
  data from outside the kernel.
\end{theorem}

\begin{proof}
  \leanok
  The generated value type is trivial, so the potential distinguishes no
  states and the step leaves every state fixed.
\end{proof}

\begin{theorem}[A one-way relation has no Rosen view]
  \label{thm:rosen-no-view}
  \lean{Systems.rosen_no_view_of_asymmetric}
  \leanok
  \uses{def:kernel}
  If some pair in a kernel's relation lacks its converse, no Rosen system
  (Rosen 1978) projects to that kernel. Rosen's relation among states is
  always symmetric.
\end{theorem}

\begin{proof}
  \leanok
  A Rosen system's relation is indistinguishability by observables, which is
  symmetric.
\end{proof}

\begin{problem}[Time and environment]
  \label{prob:time-environment}
  \uses{thm:view-roundtrip, thm:spivak-static}
  The kernel holds places for time and environment but leaves them empty.
  Find which time and environment content, if any, can be generated from the
  kernel, and state exactly what must be supplied from outside.
\end{problem}

\begin{problem}[Emergence]
  \label{prob:emergence}
  \uses{def:kernel}
  Bunge, \emph{Treatise on Basic Philosophy}, vol.~4 (1979), Definitions
  1.13 and 1.14, defines emergent properties set-theoretically. State
  emergence over the kernel without building the emergent structure into
  the model.
\end{problem}

\chapter{Protocols}

\begin{definition}[Protocol]
  \label{def:protocol}
  \uses{def:klir-system}
  A protocol is a relation $R$ over a set $T$ of roles. It exists whether or
  not anyone takes it up. This is a working definition and is not yet
  formalized.
\end{definition}

\begin{definition}[Protocol system]
  \label{def:protocol-system}
  \uses{def:protocol}
  A protocol system is a protocol together with participants acting in
  relation to it. Walch, \emph{The Fundamentals of Protocol Systems} (2023),
  p.~5: ``A protocol system is made up of a protocol, at least two people,
  and those people taking actions in relation to the protocol.'' This is a
  working definition and is not yet formalized.
\end{definition}

\begin{remark}
  A protocol is the rules, and rules are relations. By Klir, systemhood
  resides in the relation, so a protocol is the part of a system where
  systemhood lives. In computing terms, a protocol is a specification and a
  protocol system is a system running it.
\end{remark}

\begin{problem}[Formalize protocols]
  \label{prob:formalize-protocol}
  \uses{def:protocol, def:protocol-system, def:kernel}
  Define a protocol as a relation over roles and a protocol system as an
  assignment of participants to those roles, and prove that every protocol
  system is a kernel.
\end{problem}

\chapter{Three instances}

\begin{definition}[Contacts since the last wash]
  \label{def:hw-step}
  \lean{HandWashing.sanitizePeriod, HandWashing.contaminationStep}
  \leanok
  \uses{def:klir-system}
  With $N = 3$, one contact updates the count $c \in \mathbb{N}$ by
  \[
    f(c) =
    \begin{cases}
      0 & \text{if } c + 1 \ge N, \\
      c + 1 & \text{otherwise.}
    \end{cases}
  \]
\end{definition}

\begin{theorem}[Hand washing keeps contamination bounded]
  \label{thm:hand-washing}
  \lean{HandWashing.handwashing_contamination_bounded}
  \leanok
  \uses{def:hw-step}
  For all $n, c \in \mathbb{N}$, if $c \le N$ then $f^{n}(c) \le N$.
\end{theorem}

\begin{proof}
  \leanok
  Induction on $n$, using that one step never exceeds $N$.
\end{proof}

\begin{remark}
  The model is coarse: hands and the rule, against surfaces and microbes. It
  assumes the rule is followed.
\end{remark}

\begin{definition}[Organized]
  \label{def:organized}
  \lean{Systems.IsOrganized, Systems.ConcreteSystem.compose}
  \leanok
  \uses{def:klir-system}
  A set $S$ is organized when $\exists\, a, b \in S$ with $a \neq b$ and
  $\Bonded(a, b)$. Two systems with disjoint compositions and at least one
  bond between them compose by union.
\end{definition}

\begin{theorem}[The TCP and HTTP stack is organized]
  \label{thm:tcp-http}
  \lean{TCPIP.stack_organized}
  \leanok
  \uses{def:organized}
  $\Org(C_{\mathrm{HTTP}} \cup C_{\mathrm{TCP}})$, where the bond is the HTTP
  client handing data to the TCP sender.
\end{theorem}

\begin{proof}
  \leanok
  Composition of two systems with a cross bond is organized.
\end{proof}

\begin{remark}
  Each layer is modeled as closed, with no network environment. This is
  composition of two systems, not composition of arrows.
\end{remark}

\begin{definition}[Difficulty retarget as a homeostat]
  \label{def:homeostat}
  \lean{Systems.Homeostat, Systems.Homeostat.feedbackLaw, Bitcoin.difficultyHomeostat}
  \leanok
  The state is the inter-block time $s \in \mathbb{Z}$ in minutes, with set
  point $10$. The feedback law is
  \[
    \varphi(s) = s - (s - 10).
  \]
\end{definition}

\begin{theorem}[The ten-minute target is a fixed point]
  \label{thm:bitcoin}
  \lean{Bitcoin.difficulty_target_is_equilibrium}
  \leanok
  \uses{def:homeostat}
  $\varphi(10) = 10$.
\end{theorem}

\begin{proof}
  \leanok
  At the set point the error is zero and the correction leaves the state
  unchanged.
\end{proof}

\begin{remark}
  This is an idealized one-step model of the 2016-block retarget. As written,
  $\varphi$ returns $10$ from any state, so the result concerns the
  homeostat's structure, not Bitcoin's dynamics. The real rule adjusts
  difficulty multiplicatively over a 2016-block window.
\end{remark}

\chapter{Open problems}

\begin{problem}[Measuring hardness]
  \label{prob:hardness}
  Hardness can be located in a protocol by drawing its dependencies. Find a
  measure of hardness that compares protocols of different kinds.
\end{problem}

\begin{problem}[A fifth formalism]
  \label{prob:fifth-formalism}
  \uses{thm:single-arrow}
  State machines, process calculi, temporal logics, and Petri nets can be
  drawn over a common skeleton. Encode a fifth formalism and test whether the
  skeleton survives.
\end{problem}

\begin{problem}[Designed hardness]
  \label{prob:blygger}
  \uses{def:protocol}
  Blygger replaces deletion with withdrawal and makes pins irrevocable. Model
  it as an instance and locate its hardness.
\end{problem}
